\documentclass[a4paper,11pt]{article}
\usepackage{exam-style-341}
\examinehead{341 农业综合知识三 · 网络技术分章练习}{第02章 参考答案与解析}
\begin{document}

\unitanswerscover{341 农业综合知识三 · 网络技术分章练习}{第02章参考答案}{协议与体系结构：逐题答案与解析，含 OSI、TCP/IP 与五层体系结构对照。}

\answerpart{一、单项选择题（10 题，每题 3 分，共 30 分）}

\q A。\begin{analysis}网络协议由语法、语义和同步（时序）三个要素组成（教材 1.7.2，印刷页 28—29）。语法规定数据与控制信息的格式；语义规定发出何种控制信息、完成何种动作及做出何种应答；同步规定事件实现顺序。\end{analysis}

\q A。\begin{analysis}计算机网络体系结构是计算机网络各层及其协议的集合，即网络的分层结构与各层协议的统称（教材 1.7.1，印刷页 27—28）。B、C、D 分别只涉及硬件、软件或拓扑，都不完整。\end{analysis}

\q B。\begin{analysis}运输层负责向两台主机中进程之间的通信提供通用的数据传输服务，可实现端到端的可靠传输（TCP）；网络层主要提供主机到主机的路由与转发，数据链路层负责相邻节点间的帧传输（教材 1.7.3）。\end{analysis}

\q B。\begin{analysis}OSI 参考模型自下而上为物理层、数据链路层、网络层、运输层、会话层、表示层、应用层，第三层是网络层（教材 1.7.3，印刷页 29—32）。\end{analysis}

\q A。\begin{analysis}TCP/IP 体系结构的网际层使用 IP 协议，它是整个 TCP/IP 体系的核心（教材 1.7.5，印刷页 36）。TCP、UDP 属于运输层，HTTP 属于应用层。\end{analysis}

\q B。\begin{analysis}五层原理体系结构中，各层 PDU 自下而上为：比特（物理层）、帧（数据链路层）、分组／数据报（网络层）、报文段（运输层）、报文（应用层）。\end{analysis}

\q C。\begin{analysis}路由器工作在网络层，根据分组首部中的目的 IP 地址查找转发表进行转发（教材 1.7.3、4.x）。\end{analysis}

\q D。\begin{analysis}应用层的数据单元通常称为报文；运输层为报文段，网络层为分组（数据报），数据链路层为帧，物理层为比特。\end{analysis}

\q B。\begin{analysis}OSI 表示层负责数据表示、加密、压缩等与数据表示有关的功能；会话层负责建立、管理和终止会话（教材 1.7.3）。\end{analysis}

\q B。\begin{analysis}客户／服务器方式中，客户是服务请求方，主动发起请求；服务器是服务提供方，被动等待并响应请求（教材 1.3.1，印刷页 10）。\end{analysis}

\answerpart{二、多项选择题（3 题，每题 4 分，共 12 分）}

\q ABC。\begin{analysis}OSI 七层为物理层、数据链路层、网络层、运输层、会话层、表示层、应用层。D 网络接口层属于 TCP/IP 体系结构，不是 OSI 的层次。\end{analysis}

\q ABC。\begin{analysis}TCP/IP 四层为网络接口层、网际层、运输层、应用层。D 会话层属于 OSI 参考模型，TCP/IP 中没有独立的会话层（教材 1.7.5）。\end{analysis}

\q ABC。\begin{analysis}A、B 正确：发送时逐层加首部封装，接收时逐层剥去首部解封装；C 正确：物理层传输的是比特流；D 错误：封装发生在每一层，不只是运输层，例如数据链路层加帧首部和帧尾部、网络层加 IP 首部。\end{analysis}

\answerpart{三、判断题（6 题，每题 2 分，共 12 分）}

\q √。\begin{analysis}网络协议的三要素是语法、语义和同步（时序），三者缺一不可（教材 1.7.2，印刷页 28—29）。\end{analysis}

\q √。\begin{analysis}OSI 参考模型分 7 层，TCP/IP 体系结构分 4 层；五层原理体系结构综合二者优点，是教学模型（教材 1.7.3、1.7.5）。\end{analysis}

\q ×。改正：服务是下层向上层提供的功能，协议是同层对等实体之间进行通信的规则。\begin{analysis}服务是垂直方向（相邻层之间）的关系，协议是水平方向（同层实体之间）的关系；协议实现服务，上层使用下层提供的服务（教材 1.7.4，印刷页 34—35）。\end{analysis}

\q ×。改正：在 OSI 参考模型中，运输层负责端到端的可靠传输，网络层负责路由选择和分组转发。\begin{analysis}教材 1.7.3 指出运输层向两台主机中进程之间的通信提供通用的数据传输服务；网络层负责把分组从源主机送到目的主机，主要提供主机到主机的通信。\end{analysis}

\q ×。改正：数据链路层的协议数据单元是帧，运输层的协议数据单元才是报文段。\begin{analysis}各层 PDU 为：物理层比特、数据链路层帧、网络层分组、运输层报文段、应用层报文。\end{analysis}

\q √。\begin{analysis}客户／服务器方式中客户主动请求服务，服务器被动等待并响应请求，二者角色明确（教材 1.3.1，印刷页 10）。\end{analysis}

\answerpart{四、填空题（6 空，每空 1 分，共 6 分）}

\q 同步（或时序）。\begin{analysis}协议三要素为语法、语义和同步（教材 1.7.2）。\end{analysis}

\q 7。\begin{analysis}OSI 参考模型共 7 层（教材 1.7.3）。\end{analysis}

\q 应用层。\begin{analysis}TCP/IP 自下而上为网络接口层、网际层、运输层、应用层，最上层是应用层（教材 1.7.5）。\end{analysis}

\q 分组（或数据报）。\begin{analysis}网络层的 PDU 是分组，又称数据报（教材 1.7.3）。\end{analysis}

\q 帧。\begin{analysis}数据链路层把网络层交下来的分组封装成帧，其 PDU 是帧（教材 3.1.1）。\end{analysis}

\q RFC。\begin{analysis}互联网的正式标准都以 RFC 文档的形式发表，由 IETF 等组织维护（教材 1.2.3，印刷页 8）。\end{analysis}

\answerpart{五、简答题（2 题，每题 10 分，共 20 分）}

\q 参考答案要点：
（1）三要素（6 分）：语法——规定数据与控制信息的格式，例如帧中字段的排列顺序和每段长度；语义——规定需要发出何种控制信息、完成何种动作以及做出何种应答，例如确认帧的含义；同步（时序）——规定各种事件实现顺序的详细说明，例如先建立连接再传数据。
（2）分层的好处（4 分）：把复杂问题分解为若干相对独立的子问题，各层只关注本层功能；某层实现变化不影响其他层，灵活性好；各层便于实现、调试和维护；有利于标准化和不同厂商设备互操作；上层使用下层提供的服务，接口清晰。（教材 1.7.2，印刷页 27—29）

\q 参考答案要点：
（1）各层名称（4 分）：OSI 七层自下而上为物理层、数据链路层、网络层、运输层、会话层、表示层、应用层；TCP/IP 四层自下而上为网络接口层、网际层、运输层、应用层；五层为物理层、数据链路层、网络层、运输层、应用层。
（2）区别与联系（3 分）：OSI 先有模型后实现，层次划分完整、理论严谨，但过于复杂；TCP/IP 先有协议后有模型，简单实用，已成为互联网事实标准，但其网络接口层没有细分，理论体系不够完整。五层体系结构综合二者优点，是教学用的原理模型，工程上仍以 TCP/IP 为准（教材 1.7.5，印刷页 36）。
（3）五层各层主要功能（3 分）：物理层在传输媒体上传输比特流；数据链路层把分组封装成帧，实现相邻节点间的帧传输和差错检测；网络层负责路由选择和分组转发；运输层为两台主机中进程之间提供端到端的通信服务；应用层直接为用户应用进程提供服务。

\answerpart{六、综合题（1 题，共 20 分）}

\q 参考答案要点：
（1）协议与 PDU（5 分）：应用层使用 HTTP，PDU 为报文；运输层使用 TCP，PDU 为报文段（segment）；网络层使用 IP，PDU 为分组（数据报）；数据链路层使用以太网协议（802.3），PDU 为帧；物理层传输比特流，单位为比特。
（2）封装与解封装（6 分）：发送端应用层把 HTTP 报文交给运输层，运输层加 TCP 首部形成报文段；网络层加 IP 首部形成分组；数据链路层加帧首部和帧尾部形成帧；物理层把帧转换为比特流发送到传输媒体上。接收端过程相反：物理层收到比特流后交给数据链路层，数据链路层检查 FCS 并剥去帧首部／尾部，网络层剥去 IP 首部，运输层剥去 TCP 首部，最后应用层得到 HTTP 报文，实现解封装。
（3）设备与转发依据（5 分）：交换机工作在数据链路层，转发时依据帧中的目的 MAC 地址查找 MAC 地址表；路由器工作在网络层，转发时依据分组中的目的 IP 地址查找路由表（转发表）。
（4）体系结构对应关系（4 分）：OSI 应用层、表示层、会话层对应 TCP/IP 应用层和五层应用层；OSI 运输层对应 TCP/IP 运输层和五层运输层；OSI 网络层对应 TCP/IP 网际层和五层网络层；OSI 数据链路层与物理层对应 TCP/IP 网络接口层和五层的这两层。互联网实际使用的是 TCP/IP 体系结构；OSI 是理论参考模型，五层是教学模型。

\end{document}
